% volume-ii/natural-numbers/notes/arithmetic/addition/notes-landau-lean-proofs.tex
%
% Line-by-line annotation of the Lean 4 proofs of Theorem 4 (Landau).
% Purpose: to memorialize Lean arcana encountered during Project Landau
% and to serve as a reference for future Lean proof work in this project.

\subsection*{Lean 4 Proofs: Theorem 4 (Landau)}

The following section annotates the Lean 4 proofs of the two halves of
Landau's Theorem~4. Both proofs live in
\texttt{lean/LRA/VolumeII/NaturalNumbers/Addition.lean} and are verified
without Mathlib, using only the axioms in \texttt{Axioms.lean}.

% ------------------------------------------------------------------
\subsubsection*{Part I: Uniqueness (\texttt{addition\_unique})}
% ------------------------------------------------------------------

\begin{verbatim}
theorem addition_unique :
    ∀ f g : N → N → N,
    (∀ x : N, f x one = successor x) →
    (∀ x : N, g x one = successor x) →
    (∀ x y : N, f x (successor y) = successor (f x y)) →
    (∀ x y : N, g x (successor y) = successor (g x y)) →
    ∀ x y : N, f x y = g x y := by
\end{verbatim}

\noindent
The theorem statement is a chain of implications. In Lean~4, every
\verb|→| is a function arrow: the theorem is literally a function that
takes hypotheses as arguments and returns a proof. The type
\verb|N → N → N| is curried --- it means a function that takes an
\verb|N| and returns a function \verb|N → N|, which is equivalent to a
function on $\mathbb{N} \times \mathbb{N}$.
The \verb|∀ f g : N → N → N| quantifies over the two functions just as
\verb|∀ x : N| quantifies over natural numbers; functions are first-class
values in Lean's type theory.

\begin{verbatim}
    intro f g
    intro hfx
    intro hgx
    intro hfs
    intro hgs
    intro x y
\end{verbatim}

\noindent
The \verb|intro| tactic peels off one layer of \verb|∀| or \verb|→| at a
time, moving the universally quantified variable or hypothesis into the
local context. Since \verb|f| and \verb|g| appear first under \verb|∀|,
they must be introduced before the hypotheses that mention them.
The names \verb|hfx|, \verb|hgx|, \verb|hfs|, \verb|hgs| are chosen
mnemonically: \verb|hfx| is the hypothesis about \verb|f| at base,
\verb|hfs| is the hypothesis about \verb|f| at the successor step, and
so on. After \verb|intro x y|, both $x$ and $y$ are fixed arbitrary
elements of $\mathbb{N}$ and the goal is \verb|f x y = g x y|.

\begin{verbatim}
    apply A5_Induction (fun v => f x v = g x v)
\end{verbatim}

\noindent
This is the central move. \verb|A5_Induction| has type
\[
  \forall (P : \mathbb{N} \to \mathsf{Prop}),\;
  P(1) \to
  (\forall x,\, P(x) \to P(x')) \to
  \forall x,\, P(x).
\]
We instantiate $P$ as the predicate \verb|fun v => f x v = g x v|,
which says ``$f$ and $g$ agree at $v$ when the first argument is the
fixed $x$.'' Note that $x$ is free in this predicate --- it is captured
from the surrounding context. Lean creates two new goals: the base case
$P(1)$ and the inductive step.

\begin{verbatim}
    · rw [hfx x, hgx x]
\end{verbatim}

\noindent
The base case goal is \verb|f x one = g x one|.
The \verb|rw| tactic rewrites using equalities: \verb|hfx x| is the
specialisation of \verb|hfx| to the fixed \verb|x|, giving
\verb|f x one = successor x|. After rewriting, \verb|f x one| becomes
\verb|successor x|. Then \verb|hgx x| rewrites \verb|g x one| to
\verb|successor x| as well. The goal becomes
\verb|successor x = successor x|, which is closed by \verb|rfl|
automatically.

\begin{verbatim}
    · intro v ih
      rw [hfs x v, hgs x v]
      congr
\end{verbatim}

\noindent
The inductive step goal is
\verb|f x (successor v) = g x (successor v)|,
with inductive hypothesis \verb|ih : f x v = g x v| in context.
\verb|hfs x v| rewrites \verb|f x (successor v)| to
\verb|successor (f x v)|, and \verb|hgs x v| rewrites
\verb|g x (successor v)| to \verb|successor (g x v)|.
The goal is now \verb|successor (f x v) = successor (g x v)|.
The \verb|congr| tactic (``congruence'') strips the outermost
constructor --- here \verb|successor| --- reducing the goal to
\verb|f x v = g x v|, which is exactly \verb|ih|. Lean closes
the goal automatically.

% ------------------------------------------------------------------
\subsubsection*{Part II: Existence (\texttt{addition\_exists})}
% ------------------------------------------------------------------

\begin{verbatim}
theorem addition_exists :
    ∃ f : N → N → N,
    (∀ x : N, f x one = successor x) ∧
    (∀ x y : N, f x (successor y) = successor (f x y)) := by
\end{verbatim}

\noindent
An existential goal \verb|∃ f, P f| requires producing a witness ---
a concrete \verb|f| --- and a proof that \verb|P f| holds.

\begin{verbatim}
  let f : N → N → N :=
    fun x y =>
      N_rec (successor x) (fun prev => successor prev) y
\end{verbatim}

\noindent
The \verb|let| tactic introduces a local definition into the proof
context. The witness \verb|f| is constructed using \verb|N_rec|, the
primitive recursion axiom on $\mathbb{N}$:
\[
  \mathtt{N\_rec} : T \to (T \to T) \to \mathbb{N} \to T
\]
Applied here with $T = \mathbb{N}$:
\begin{itemize}
  \item Base value: \verb|successor x| --- the result when $y = 1$.
  \item Step function: \verb|fun prev => successor prev| --- at each
        successor step, wrap the previous result in \verb|successor|.
  \item Argument: \verb|y| --- the natural number to recurse on.
\end{itemize}
This constructs exactly the piecewise definition
\[
  f(x, y) = \begin{cases} x' & \text{if } y = 1 \\ (f(x,v))' & \text{if } y = v' \end{cases}
\]

\begin{verbatim}
  refine ⟨f, ?_, ?_⟩
\end{verbatim}

\noindent
The \verb|refine| tactic provides a partial proof term with holes marked
\verb|?_|. The angle bracket notation \verb|⟨f, ?_, ?_⟩| constructs the
existential witness: \verb|f| is the function, and the two \verb|?_|
placeholders become the two conjuncts to prove. Lean creates two new
goals, one for each.

\begin{verbatim}
  · intro x
    show N_rec (successor x) (fun prev => successor prev) one = successor x
    rw [N_rec_one]
\end{verbatim}

\noindent
The first goal is \verb|f x one = successor x|. Because \verb|f| is a
\verb|let|-binding, Lean does not automatically unfold it. The
\verb|show| tactic replaces the goal with a definitionally equal
expression --- here the explicit unfolding of \verb|f x one| to its
\verb|N_rec| form --- making the pattern visible for rewriting.
\verb|N_rec_one| is the computation rule
\verb|N_rec base step one = base|, so the rewrite closes the goal.

\begin{verbatim}
  · intro x y
    show N_rec (successor x) (fun prev => successor prev) (successor y) =
         successor (N_rec (successor x) (fun prev => successor prev) y)
    rw [N_rec_succ]
\end{verbatim}

\noindent
The second goal is \verb|f x (successor y) = successor (f x y)|.
Again \verb|show| unfolds \verb|f| explicitly on both sides.
\verb|N_rec_succ| is the computation rule
\verb|N_rec base step (successor n) = step (N_rec base step n)|.
After rewriting, the step function \verb|fun prev => successor prev|
applied to \verb|N_rec ... y| reduces to
\verb|successor (N_rec ... y)|, which is exactly the right-hand side.
The goal is closed.

\begin{remark*}[The \texttt{show} tactic and definitional equality]
In Lean~4, \verb|let|-bound definitions are definitionally equal to
their bodies but are not automatically unfolded by tactics like
\verb|rw|. The \verb|show| tactic is the standard way to replace a goal
with a definitionally equal expression, forcing the unfolding and making
syntactic patterns available for rewriting. This is a recurring pattern
when working with locally constructed witnesses.
\end{remark*}

\begin{remark*}[The \texttt{congr} tactic]
The \verb|congr| tactic proves goals of the form $f\,a = f\,b$ by
reducing them to $a = b$, exploiting the fact that functions respect
equality (congruence). It is particularly useful at the end of inductive
steps where the outermost constructor is known to match and only the
recursive arguments need to be compared.
\end{remark*}

\begin{remark*}[Lean 4 verification]
Both theorems are verified in
\texttt{lean/LRA/VolumeII/NaturalNumbers/Addition.lean}
without Mathlib, using only \texttt{A5\_Induction}, \texttt{N\_rec},
\texttt{N\_rec\_one}, and \texttt{N\_rec\_succ} from \texttt{Axioms.lean}.
\end{remark*}